\scanpage{152}
\[
\unsure{X|\mathcal X|\mathfrak X}_{P\mathcal D^+}\simeq PD^+\times\mathbb C/\mathbb Z^2\tag*{(20)}
\]
où $\mathbb Z^2$ opère sur $PD^+\times\mathbb C$ par
\[
T_{m,n}(z_1,z_2,z)=(z_1,z_2,z+m\deleted{\unsure{Z|z}}z_1+nz_2)
\]

Alors le groupe $\mathbb C^*\times SL(2,\mathbb Z)$ opère sur $X_{PD^+}\longrightarrow PD^+$, en faisant opérer $SL(2,\mathbb Z)$ par
\[
T_g(z_1,z_2,z)=(dz_1+cz_2,bz_1+az_2,z)\tag*{(21)}
\]
et $\mathbb C^*$ \deleted{\unsure{$T$|$\lambda$} \illegible{2} \unsure{$(20)$}} par
\[
\deleted{\unsure{\lambda|T_m|T_n}}\ T_\lambda(z_1,z_2,z)=\left(\frac{1}{\lambda}z_1,\frac{1}{\lambda}z_2,\frac{1}{\lambda}z\right).\tag*{(22)}
\]
% slanted
\margin{opération de $\mathbb C^*$ définie par le pb modulaire, la base $\varphi$ \inserted{(de $t_E$)} multipliée par $\lambda\varphi$ \ldots}

Chacun des deux groupes \unsure{commutants|commutant|commute} $\mathbb C^*,SL(2,\mathbb Z)$ opère librement sur $X_{PD^+}$ et même sur $PD^+$, mais l'opération du groupe produit $\mathbb C^*\times SL(2,\mathbb Z)$ \inserted{sur $PD^+$} n'est pas une opération libre, p.ex. l'élément $(-1,-1)$ opère trivialement sur $PD^+$ (autom. universel des courbes elliptiques) -- mais non sur $X_{PD^+}$, mais il \deleted{\unsure{opère|n'a pas}} a des pts fixes \deleted{\unsure{à savoir|de moi}} $X_{PD^+}$ (et d'autres ?!)

\deleted{En fait, l'opération sur $X_{PD^+}$ est une opération libre, sauf en les \unsure{T|I}}

\deleted{quotient par $\mathbb C^*$ étant $X_{D^+}$, et le quotient de ce dernier par $SL(2,\mathbb Z)$, \inserted{moins la section nulle} s'identifiant à $M_{1,2}$ !}
% archive sheet number: 151
